\section{A Limit for Every Protocol}
\label{sec:limit}

The first result applies to every protocol in the model, not only to \sys{}. It uses one
property of verifiers: they are deterministic, and their only inputs are messages and
clock readings.

\begin{prp}[Isolation reach implies hostile tenure]
\label{prp:limit}
Let $f \ge 1$. If a protocol has isolation reach $P$ at a verifier $v$, then it has hostile
tenure $P$ at $v$. In the hostile run, $v$ accepts the same commands at the same times as
in the isolated run.
\end{prp}

\begin{argument}[Candidate argument (sketch)]
Take the isolated run $\rho$ of \cref{def:reach}. Build $\rho'$ equal to $\rho$ until
$t_0$. At $t_0$ the adversary obtains the key and state of the holder $c$, and every other
officer withdraws. From $t_0$ the adversary drops every message from the other parties to
$c$ and $v$. It runs a copy of $c$ that makes the choices $c$ made in $\rho$, and it
delivers to $v$ exactly the messages that $v$ received in $\rho$, at the same times. The
clock of $v$ is the same function of real time in both runs. So $v$ has the same inputs,
and it accepts the same commands at the same times. \Cref{app:arguments} gives the steps.
\end{argument}

This is the indistinguishability argument of CAP-style
results~\cite{gilbert2002cap}, and the trade-off it expresses is the one that leases
make explicit~\cite{gray1989leases}. We restate it in terms of isolation reach and
withdrawal and claim nothing new about the technique \claim{C-LIM}. Four remarks make the
statement precise.
\begin{itemize}
  \item It is existential. It says that some run has hostile tenure $P$, not that service
  continues in every isolated run (that is \cref{def:avail}).
  \item The proof never uses the fact that $P$ is a time. A limit on the number of
  commands has the same property.
  \item The verifier accepts the same commands in both runs. Checks of command content
  therefore cannot shorten the tenure. They can only reject commands that a loyal holder
  would not send either.
  \item An input that the adversary does not control, such as a person who sets a
  local-control switch at the device, is outside the argument.
\end{itemize}
A designer who wants zero hostile tenure must therefore accept zero isolation reach, as
with a quorum signature on every command. Any positive reach is also a minimum hostile
tenure.
